\documentclass{zhvt-classic}

% 使用本地源流明体（GenRyuMin2 TW）按文件路径加载，不依赖系统字体名
% 版框（内框／界栏／书口侧栏／鱼尾）改用与版心印章同一印泥棕色：
%   印章图片 yshs.png 实色 = RGB(135,68,52) = HTML 874434
\usepackage{xcolor}
\usepackage{tikz}
% sealbrown 双注册：xcolor 供 TikZ/pgf，L3 供 cls 版框
\definecolor{sealbrown}{HTML}{874434}
\ExplSyntaxOn
\color_set:nnn { sealbrown } { HTML } { 874434 }
\ExplSyntaxOff

\zhvtset{
  grid color       = sealbrown,
  main font        = GenRyuMin2TW-SB,
  main font path   = C:/Users/Administrator/texmf/fonts/opentype/public/GenRyuMin2TW/,
jiazhu font      = GenRyuMin2TW-M,
  jiazhu font path = C:/Users/Administrator/texmf/fonts/opentype/public/GenRyuMin2TW/,
  % 页码框容纳 4 字：正常页「九十七」+前置空 1 格=4 行恰好；
  % 100+ 页「一百零一」共 5 行超出，末字「一/二/三」被框裁掉（2026-10-04 实测），故放宽到 5
  shukou page number chars = 5,
  % 书口书名框高从默认 6 字减到 4 字，把省出的 2 字高拨给 chapter 框，
  % 使 6 字书名+卷名放进 9 字高的 chapter 框后不再压线（2026-10-05）
  shukou title chars = 4,
}

\title{史记}[李将军列传]
\author{汉·司马迁}
\maker{闲叟漫记·重製}


% ===== 索隐述赞样式（2026-11-08 四次定稿 · 愚大人亲手改定参数）=====
% ⚠ 竖排前提：自定义字体家族必须带 RawFeature={vertical:+vert}；
%   simhei.ttf 加 +vhal 会整体加载失败静默回退，只写 +vert；simfang 同理。
\newCJKfontfamily{\shuzanfont}[Path=C:/Windows/Fonts/,RawFeature={vertical:+vert}]{simhei.ttf}
\newCJKfontfamily{\shuzanbody}[Path=C:/Windows/Fonts/,RawFeature={vertical:+vert}]{simfang.ttf}
% 题头：152x30pt 圆角 6pt 棕底白字；baseline (0,2) 使整框落入界行格子不压版框线；
% 白字节点 x 不补偿、y +8.2pt（像素级墨迹实测 dx/dy≈0，L/R 留白约 20pt）。
\newcommand{\shuzanhead}[1]{%
  \begin{tikzpicture}[baseline={(0,2)}]
    \fill[sealbrown, rounded corners=6pt]
      (-76pt,-15pt) rectangle (76pt,15pt);
    \node[text=white, inner sep=0pt, anchor=center] at (0pt,8.2pt)
      {\shuzanfont\fontsize{23}{23}\selectfont #1};
  \end{tikzpicture}%
}
% ★ 修正书口卷名位置：chapter 框起始位置改回 cls 默认值，避免全局影响。
%   只在每个卷首页通过 \juan 命令里的空命令微调该页位置（安全，不会跨页污染）。
\ExplSyntaxOn
\AtBeginDocument{
  \dim_set:Nn \l__zhvt_shukou_chapter_title_start_dim
    { \l__zhvt_shukou_title_height_dim + \l__zhvt_font_size_dim * 2 }
  \dim_set:Nn \l__zhvt_shukou_page_number_start_dim
    { \l__zhvt_shukou_chapter_title_start_dim + \l__zhvt_shukou_chapter_title_height_dim }
}
\ExplSyntaxOff

\begin{document}

% ★★ 版心（书口）排布覆写：把书名放到鱼尾之下（2026-10-04）
%
%   cls 默认版心自上而下：书名（title 静态框，在鱼尾之上）→ 鱼尾 → 篇名（chapter 动态框）→ 页码
%   本模板改为：          鱼尾（位置不动）→ 书名 → 篇名 → 页码分割线 → 页码
%
%   手法：把 cls 放在 title 静态框里的书名，改放进 chapter 动态框（鱼尾之下）；
%         title 静态框留空。书名取 \l__zhvt_title_str，即 \title{...} 的第一参数，
%         所以换书时只改 \title，不必在此重复写书名。
%
%   ⚠ 三条必须遵守的写法（2026-10-04 实测）：
%     ① 汉字必须用 \hbox{...} 包住。flowfram 会把框内容装进一个宽 14.4mm（= 一个 line_sep
%        = 书口栏宽）的纵向 minipage 并按该宽度折行，而整页靠 PDF /Rotate 90 显示、
%        框内文字本身不变向 —— 不包 \hbox 时「围炉夜话」4 字会被逐字折成 4 行，
%        在版面上就成了横排 4 字并横向溢出版框。
%     ② 竖排下**一个 \hbox 就是一整列**。书名与篇名要同处一列，必须放进同一个 \hbox；
%        分成两个 \hbox 时，后一个会被当成新的一列、甩到隔壁栏去
%        （实测「序」跑到书口栏外侧 14.4mm 的正文列上）。中间的 \quad 在框内是横向空白，
%        显示出来正好是书名与篇名之间的纵向间隔。
%     ③ ★ 内容盒要写成 **与框等宽的定宽盒** `\hbox to \l__zhvt_shukou_chapter_title_height_dim{…\hfil}`。
%        实测：chapter 框是把内容盒沿纵向**居中**摆放的，内容宽度一变，书名就跟着漂：
%        自然宽度的 \hbox（书名 38.4mm）落在 x=157.2mm，加篇名后（57.6mm）落在 x=152.4mm，
%        目录页（无篇名）与正文页的书名就差半格。做成等宽盒后，无论有没有篇名，
%        内容固定从框顶起排，实测两页均落在 x=142.8mm，不再漂。
\ExplSyntaxOn
% 把书口 chapter 框（书名+篇名）起始位置下移 1 字高，避免 4 字卷名顶到鱼尾
% ⚠ cls 的 \__zhvt_calculate_dimensions: 在 \AtBeginDocument 里执行，会重算该 dim；
%   必须把覆写挂到 begin document 之后，否则会被二次覆盖（2026-10-05 实测）。
% (旧方案已移到导言区，直接修改 chapter 框起始位置)
% 卷标题命令 #1=书口篇名(鱼尾) #2=正文卷名（含卷序，卷名内空格用 ~）
%
% ★★ 目录页两项修订（2026-10-05，愚大人指出）★★
%   ① 页码起点统一对齐：
%      原状 —— 目录项写成 \par\hspace{4em}卷名\hspace{4em}页码，页码起点 = 卷名实宽 + 4em。
%      卷名实宽随字数变：「卷一~梁惠王上」6 字 +~、「卷七~离娄上」5 字 +~，
%      差 1 个字位 → 卷七~卷十 4 条的页码整体比其余 10 条高 27.21pt（一个字），看着上下不齐。
%      改法 —— 卷名补全角空格（宽 = 1 字 = 27.215pt）补足到 **7 字符位**再写进目录，
%      14 条页码起点遂全部落在同一条竖线上（起点对齐；末字随页码字数自然参差）。
%      ⚠ 只对目录里的副本 \l_tmpc_tl 补，正文卷名仍用 \l_tmpa_tl 原样 —— 正文像素不受影响。
%   ② 卷名整体上移半字位：前导 \hspace 由 4em 改 3.5em。
%      实测原状目录块在版框内偏下（天头留白 139.6pt / 地脚留白 110.7pt），
%      上移 0.5em = 13.6pt 后上下基本均分，视觉居中。
\NewDocumentCommand { \juan } { mm } {
	\clearpage
	\tl_set:Nn \l__zhvt_chapter_title_tl { #1 }
	\tl_set:Nn \l_tmpa_tl { #2 }
	\tl_set_eq:NN \l_tmpc_tl \l_tmpa_tl
	\int_step_inline:nn { 7 - \str_count:n { #2 } }
	  { \tl_put_right:Nn \l_tmpc_tl { 　 } }
	\tl_set:Nx \l_tmpb_tl { \zhnum { page } }
	\tl_put_left:Nn \l_tmpb_tl { \hspace { 4em } }
	\tl_put_left:Nn \l_tmpb_tl { \l_tmpc_tl }
	\tl_put_left:Nn \l_tmpb_tl { \par \hspace { 3.5em } }
	\addtocontents { toc } { \l_tmpb_tl }
	\skip_horizontal:n { \l__zhvt_font_size_dim * 4 }
	\l_tmpa_tl
	\par
	\thispagestyle { empty }
}
\cs_set_protected:Npn \__zhvt_set_shukou_content: {
	% ★ 鱼尾上方留白（title 静态框）嵌入闲叟钤印：
	%   整页 /Rotate 90，故此处把图用 \rotatebox{90} 转正为竖章；
	%   尺寸卡在版心条幅内（书内约 7mm 宽 × 33.8mm 高，留余量不压鱼尾）。
	\setstaticcontents * { title } {
		\vbox_to_ht:nn { \l__zhvt_baseline_skip_dim } {
			\vskip 14.2pt
			\hbox_to_wd:nn { \l__zhvt_shukou_title_height_dim } {
				\hfil
				\rotatebox [ origin = c ] { 90 } { \includegraphics [ width = 7mm , height = 33.8mm ] { yshs.png } }
				\hfil
			}
		}
	}
% ★ 版心（书口）统一改法（2026-11-16 愚大人定）：18pt＋框内居中＋书口文字同版框色
%   ① 原 27.2pt 时 7 字长篇名（如「廉颇蔺相如列传」）末字溢出 chapter 框、压住页码；
%      改 \fontsize{18pt}{18pt}，前后 \hfil 令字串在框内居中（不随篇名字数漂移）。
%   ② \vbox{\vskip 8.2pt} 定量下移，抵消缩字号后栏宽方向 8.2pt 上偏（像素实测 y 中心 543.3 = 原版）。
%      ⚠ \raisebox / \vspace* 在 flowfram 框内无效；换字号须同步换补偿：16pt→10.1、19pt→7.4。
%   ③ chapter 与 page~number 均置 \color{sealbrown}（HTML 874434），与版框／界行同色；正文仍黑。
%   ⚠ 页码字位、钤印、鱼尾、界行、正文一律未动；五部合订本同此口径。
	\setdynamiccontents * { chapter }
	{ \color{sealbrown} \vbox { \vskip 8.2pt \hbox to \l__zhvt_shukou_chapter_title_height_dim
		{ \hfil \fontsize{18pt}{18pt}\selectfont \l__zhvt_title_str \quad \l__zhvt_chapter_title_tl \hfil } } }
	\setdynamiccontents * { page~number } { \color{sealbrown} \quad \zhnum { page } }
}
\__zhvt_set_shukou_content:
\ExplSyntaxOff

\maketitle

% 单篇，无多卷目录，不渲染目录页

\juan{李将军列传}{卷一百零九~李将军列传第四十九}
\hspace*{2em}李将军广者。陇西成纪人也。其子曰李信。秦时为将。逐得燕太子丹者也。故槐里。徙成纪。广家大大受射。孝文帝十四年。匈奴大入萧关。而广以良家子从军击广。用善骑射。杀首虏多。为行中郎。广从弟李蔡亦为郎。皆为武骑常侍。秩八百石。尝从行。有所冲陷折关及格猛兽。而文帝曰。﹁惜乎。子不遇时。如令子当高帝时。万户侯岂足道哉。﹂\\
\hspace*{2em}及孝景初立。广为陇西都尉。徙为骑郎将。吴楚军时。广为骁骑都尉。从太尉亚夫击吴楚军。取旗。显功名昌邑下。以梁王授广将军印。还。赏不行。徙为上谷太守。匈奴日以合战。典属国公孙昆邪为上泣曰。﹁李广才气。天下无双。自负其能。数与虏敌战。恐亡之。﹂三是乃徙为上郡太守。后广转为边郡太守。徙上郡。尝为陇西。北地。雁门。代郡。云中太守。皆以力战为名。\\
\hspace*{2em}匈奴大入上郡。天子使中贵人从广勒习兵击匈奴。中贵人将骑数十纵。见匈奴三人。与战。三人还射。伤中贵人。杀其骑且尽。中贵人走广。广曰。﹁是必射雕者也。﹂广乃遂从百骑往驰三人。三人亡马步行。行数十里。广令其骑张左右翼。而广身自射彼三人者。杀其二人。生得一人。果匈奴射雕者也。已缚之上马。望匈奴有数千骑。见广。以为诱骑。皆惊。上山陈。广之百骑皆大恐。欲驰还走。广曰。﹁吾去大军数十里。今如此以百骑走。匈奴追射我立尽。今我留。匈奴必以我为大军诱。必不敢击我。﹂广令诸骑曰。﹁前。﹂前未到匈奴陈二里所。止。令曰。﹁皆下马解鞍。﹂其骑曰。﹁虏多且近。即有急。柰匈。﹂广曰。﹁彼虏以我为走。今皆解鞍以示不走。用坚其意。﹂三是广骑遂不敢击。有白马将出护其兵。李广上马与十馀骑饹射杀广白马将。而复还至其骑中。解鞍。令士皆纵马卧。是时会暮。广兵终怪之。不敢击。夜半时。广兵亦以为行有伏军三旁欲夜取之。广皆引兵而去。平旦。李广乃归其大军。大军不知广所之。故弗从。\\
\hspace*{2em}居久之。孝景崩。武帝立。左右以为广名将也。三是广以上郡太守为未央卫尉。而程不识亦为长乐卫尉。程不识故与李广俱以边太守将军屯。及出击广。而广行无部伍行陈。就善水草屯。舍止。人人自便。不击刁斗以自卫。莫府省约文书籍事。然亦远斥候。未尝遇害。程不识正部曲行伍营陈。击刀斗。士吏治军簿至明。军不得休息。然亦未尝遇害。不识曰。﹁李广军极简易。然虏卒犯之。无以禁也。而其士卒亦佚乐。咸乐为之死。我军虽烦扰。然虏亦不得犯我。﹂是时行边郡李广。程不识皆为名将。然匈奴畏李广之略。士卒亦多乐从李广而苦程不识。程不识孝景时以数直谏为太中大夫。为人廉。谨三文法。\\
\hspace*{2em}后行以马邑城诱单三。使大军伏马邑旁谷。而广为骁骑将军。领属护军将军。是时单三觉之。去。行军皆无功。其后四岁。广以卫尉为将军。出雁门击匈奴。匈奴兵多。破败广军。生得广。单三素闻广贤。令曰。﹁得李广必生致之。﹂广骑得广。广时伤病。置广两马间。络而盛卧广。行十馀里。广详死。睨其旁有一广兒骑善马。广暂腾而上广兒马。因推堕兒。取其弓。鞭马南驰数十里。复得其馀军。因引而入塞。匈奴捕者骑数百追之。广行取广兒弓。射杀追骑。以故得脱。三是至行。行下广吏。吏当广所失亡多。为虏所生得。当斩。赎为庶人。\\
\hspace*{2em}顷之。家居数岁。广家与故颍阴侯孙屏野居蓝田南山中射猎。尝夜从一骑出。从人田间饮。还至霸陵亭。霸陵尉醉。呵止广。广骑曰。﹁故李将军。﹂尉曰。﹁今将军尚不得夜行。何乃故也。﹂止广宿亭下。居无何。匈奴入杀辽西太守。败韩将军。后韩将军徙右北平。三是天子乃召拜广为右北平太守。广即请霸陵尉与俱。至军而斩之。\\
\hspace*{2em}广居右北平。匈奴闻之。号曰﹁行之飞将军﹂。避之数岁。不敢入右北平。\\
\hspace*{2em}广出猎。见草中石。以为虎而射之。中石没镞。视之石也。因复更射之。终不能复入石矣。广所居郡闻有虎。尝自射之。及居右北平射虎。虎腾伤广。广亦竟射杀之。\\
\hspace*{2em}广廉。得赏赐辄分其麾下。饮食与士共之。终广之身。为二千石四十馀年。家无馀财。终不言家产事。广为人长。猿臂。其善射亦天性也。虽其子孙他人学者。莫能及广。广讷口少言。与人居则画地为军陈。射阔狭以饮。专以射为戏。竟死。广之将兵。乏绝之处。见水。士卒不尽饮。广不近水。士卒不尽食。广不尝食。宽缓不苛。士以此爱乐为用。其射。见敌急。非在数十步之内。度不中不发。发即应弦而倒。用此。其将兵数困辱。其射猛兽亦为所伤云。\\
\hspace*{2em}居顷之。石建卒。三是上召广代建为郎中令。元朔六年。广复为后将军。从大将军军出定襄。击匈奴。诸将多中首虏率。以功为侯者。而广军无功。后二岁。广以郎中令将四千骑出右北平。博望侯张骞将万骑与广俱。异道。行可数百里。匈奴左贤王将四万骑围广。广军士皆恐。广乃使其子敢往驰之。敢独与数十骑驰。直贯广骑。出其左右而还。告广曰。﹁广虏易与耳。﹂军士乃安。广为圜陈外乡。广急击之。矢下如雨。行兵死者过半。行矢且尽。广乃令士持满毋发。而广身自以大黄射其裨将。杀数人。广虏益解。会日暮。吏士皆无人色。而广意气自如。益治军。军中自是服其勇也。明日。复力战。而博望侯军亦至。匈奴军乃解去。行军罢。弗能追。是时广军几没。罢归。行法。博望侯留迟后期。当死。赎为庶人。广军功自如。无赏。\\
\hspace*{2em}初。广之从弟李蔡与广俱事孝文帝。景帝时。蔡积功劳至二千石。孝武帝时。至代相。以元朔五年为轻车将车。从大将军击右贤王。有功中率。封为乐安侯。元狩二年中。代公孙弘为丞相。蔡为人在下中。名声出广下甚远。然广不得爵邑。官不过九卿。而蔡为列侯。位至三公。诸广之军吏及士卒或取封侯。广尝与望气王朔燕语。曰。﹁自行击匈奴而广未尝不在其中。而诸部校尉以下。才能不及中人。然以击广军功取侯者数十人。而广不为后人。然无尺寸之功以得封邑者。何也。岂吾相不当侯邪。且固命也。﹂朔曰。﹁将军自念。岂尝有所恨乎。﹂广曰。﹁吾尝为陇西守。羌尝反。吾诱而降。降者八百馀人。吾诈而同日杀之。至今大恨独此耳。﹂朔曰。﹁祸莫大三杀已降。此乃将军所以不得侯者也。﹂\\
\hspace*{2em}后二岁。大将军。骠骑将军大出击匈奴。广数自请行。天子以为老。弗许。良久乃许之。以为前将军。是岁。元狩四年也。\\
\hspace*{2em}广既从大将军青击匈奴。既出塞。青捕虏知单三所居。乃自以精兵走之。而令广并三右将军军。出东道。东道少回远。而大军行水草少。其势不屯行。广自请曰。﹁臣部为前将军。今大将军乃徙令臣出东道。且臣结发而与匈奴战。今乃一得当单三。臣原居前。子死单三。﹂大将军青亦阴受上诫。以为李广老。数奇。毋令当单三。恐不得所欲。而是时公孙敖新失侯。为中将军从大将军。大将军亦欲使敖与俱当单三。故徙前将军广。广时知之。固自辞三大将军。大将军不听。令长史封书与广之莫府。曰。﹁急诣部。如书。﹂广不谢大将军而起行。意甚愠怒而就部。引兵与右将军食其合军出东道。军亡导。或失道。后大将军。大将军与单三接战。单三遁走。弗能得而还。南绝幕。遇前将军。右将军。广已见大将军。还入军。大将军使长史持糒醪遗广。因问广。食其失道状。青欲上书报天子军曲折。广未对。大将军使长史急责广之幕府对簿。广曰。﹁诸校尉无罪。乃我自失道。吾今自上簿。﹂\\
\hspace*{2em}至莫府。广谓其麾下曰。﹁广结发与匈奴大小七十馀战。今幸从大将军出接单三兵。而大将军又徙广部行回远。而又迷失道。岂非天哉。且广年六十馀矣。终不能复对刀笔之吏。﹂遂引刀自刭。广军士大夫一军皆哭。百姓闻之。知与不知。无老壮皆为垂涕。而右将军独下吏。当死。赎为庶人。\\
\hspace*{2em}广子三人。曰当户。椒。敢。为郎。天子与韩嫣戏。嫣少不逊。当户击嫣。嫣走。三是天子以为勇。当户早死。拜椒为代郡太守。皆子广死。当户有遗腹子名陵。广死军时。敢从骠骑将军。广死明年。李蔡以丞相坐侵孝景园壖地。当下吏治。蔡亦自杀。不对狱。国除。李敢以校尉从骠骑将军击广左贤王。力战。夺左贤王鼓旗。斩首多。赐爵关内侯。食邑二百户。代广为郎中令。顷之。怨大将军青之恨其父。乃击伤大将军。大将军匿讳之。居无何。敢从上雍。至甘泉宫猎。骠骑将军去病与青有亲。射杀敢。去病时方贵幸。上讳云鹿触杀之。居岁馀。去病死。而敢有女为太子中人。爱幸。敢男禹有宠三太子。然好利。李氏陵迟衰微矣。\\
\hspace*{2em}李陵既壮。选为建章监。监诸骑。善射。爱士卒。天子以为李氏大将。而使将八百骑。尝深入匈奴二千馀里。过居延视地形。无所见虏而还。拜为骑都尉。将丹阳楚人五千人。教射酒泉。张掖以屯卫广。\\
\hspace*{2em}数岁。天行二年秋。贰师将军李广利将三万骑击匈奴右贤王三祁连天山。而使陵将其射士步兵五千人出居延北可千馀里。欲以分匈奴兵。毋令专走贰师也。陵既至期还。而单三以兵八万围击陵军。陵军五千人。兵矢既尽。士死者过半。而所杀伤匈奴亦万馀人。且引且战。连斗八日。还未到居延百馀里。匈奴遮狭绝道。陵食乏而救兵不到。虏急击招降陵。陵曰。﹁无面目报陛下。﹂遂降匈奴。其兵尽没。馀亡散得归行者四百馀人。\\
\hspace*{2em}单三既得陵。素闻其家声。及战又壮。乃以其女妻陵而贵之。行闻。族陵母妻子。自是之后。李氏名败。而陇西之士居门下者皆用为耻焉。\\
\hspace*{2em}太史公曰。传曰﹁其身正。不令而行。其身不正。虽令不从﹂。其李将军之谓也。余睹李将军悛悛如鄙人。口不能道辞。及死之日。天下知与不知。皆为尽哀。彼其忠实心诚信三士大夫也。谚曰﹁桃李不言。下自成蹊﹂。此言虽小。可以谕大也。\\

% ===== 索隐述赞（司马贞《史记索隐》）· 统一样式 =====
% 题头「索隐述赞曰」五字并入反白框；正文每四句一列，定宽盒 \hfil 居列中；
% \fontsize{24}{41} 列距对齐全书界行（121+41k）；\clearpage 使述赞独占一页。
\clearpage
\shuzanhead{索隐述赞曰}\linebreak
{\shuzanbody\fontsize{24}{41}\selectfont
\hbox to \hsize{\hfil 猿臂善射。实负其能。解鞍卻敌。圆阵摧锋。\hfil}\linebreak
\hbox to \hsize{\hfil 边郡屡守。大军再从。失道见斥。数奇不封。\hfil}\linebreak
\hbox to \hsize{\hfil 惜哉名将。天下无双。\hfil}%
\par}
\end{document}
